
\SetTituloActual{Bebidas probióticas de col y café}

\Inicio{REVIEW}{Valorización de la col y los residuos de café en bebidas fermentadas probióticas de origen vegetal}{Gael Zamudio Sánchez, Patricia Rebeca Morales Juárez, Ruben Osvaldo Torres Pineda, Diego Alejandro Ortega Peraza}

\Resumen{El presente trabajo analiza la valorización de la col y residuos de café en el desarrollo de bebidas fermentadas probióticas de origen vegetal mediante una revisión de evidencia científica. Se evaluaron investigaciones con enfoque a la viabilidad microbiana en cepas de la bacteria ácido láctica Lactobacillus plantarum donde los resultados mostraron como la col constituye una matriz adecuada para la fermentación debido a su contenido de azúcares fermentables y microbiota natural, alzando concentraciones de entre 10$^{7}$ y 10$^{13}$ UFC/mL los cuales son aptos para entrar en el umbral de la dosis necesaria para su funcionalidad como probiótico.  El meta-análisis muestra un efecto combinado significativo (p = 0.033) con una media global de 9.49$\pm$0.63 log UFC y una alta heterogeneidad entre los estudios (I$^{2}$ $\approx$ 79\%), atribuida a diferencias en sustratos y condiciones experimentales, el conjunto de datos recopilados sugiere que la combinación de col y residuos de café podría llegar a favorecer sistemas simbióticos en el desarrollo de bebidas funcionales. Sin embargo , aún se requieren estudios experimentales para estandarizar  condiciones de fermentación y evaluar su viabilidad.}{Col, Residuos de café, \textit{Lactobacillus plantarum}, fermentación}


    \begin{multicols}{2}


\section{Introducción}

En los últimos años, el interés por el consumo de alimentos funcionales ha incrementado debido a su relación con la salud intestinal y el bienestar general. Dentro de estos, los probióticos han cobrado relevancia por su capacidad de contribuir al equilibrio de la microbiota intestinal y mejorar procesos digestivos. Estos microorganismos vivos, cuando se administran en cantidades adecuadas, pueden aportar beneficios a la salud del huésped, siendo recomendadas concentraciones mínimas del orden de 10$^{8}$ unidades formadoras de colonias (UFC) para lograr efectos significativos \citep{binda2020}.

Diversos factores pueden alterar la microbiota intestinal, entre ellos el uso de antibióticos o dietas inadecuadas, lo que puede provocar síntomas como distensión abdominal, diarrea, estreñimiento o síndrome de intestino irritable \citep{suez2019}. En este contexto, el consumo de productos con propiedades probióticas se ha propuesto como una alternativa para restaurar el equilibrio intestinal y mejorar la digestión.

Sin embargo, la mayoría de los productos probióticos disponibles en el mercado son de origen lácteo, lo cual representa una limitación para una parte importante de la población, considerando que aproximadamente el 65\% de las personas presenta algún grado de intolerancia a la lactosa \citep{storhaug2017}. Aunque algunos productos fermentados contienen enzimas como la $\beta$-galactosidasa que contribuyen a la degradación de la lactosa, la presencia de compuestos residuales puede generar efectos adversos en ciertos consumidores.

Ante esta problemática, las bebidas probióticas de origen vegetal han surgido como una alternativa viable, ya que permiten evitar compuestos potencialmente irritantes y amplían el acceso a estos productos para personas con restricciones dietéticas o preferencias alimentarias específicas. Además, la creciente tendencia hacia el consumo de alimentos que favorecen la salud intestinal refleja un interés global por mantener un microbioma equilibrado, lo que ha impulsado la búsqueda de nuevas fuentes de probióticos accesibles y funcionales.

En este sentido, según \citet{parrahuertas2010}, algunos vegetales como la col representan una matriz adecuada para procesos de fermentación, debido a la presencia natural de bacterias ácido-lácticas, principalmente \textit{Lactobacillus plantarum}. Este microorganismo posee la capacidad de fermentar azúcares y producir ácido láctico, favoreciendo tanto la conservación del alimento como el desarrollo de propiedades funcionales. Estudios previos como el de \citet{yoon2006}, han demostrado que durante la fermentación de la col pueden alcanzarse concentraciones cercanas a 10$^{9}$ UFC/mL, lo que sugiere su potencial como fuente de probióticos viables.

Por otro lado, el aprovechamiento de subproductos agroindustriales ha cobrado relevancia en el desarrollo de estrategias sostenibles. En particular, los residuos de café contienen compuestos bioactivos como melanoidinas y oligosacáridos con posible efecto prebiótico, los cuales pueden favorecer el crecimiento de bacterias benéficas y potenciar el valor funcional de los productos fermentados \citep{esquiveljimenez2012}. Su reutilización no solo contribuye a la reducción del impacto ambiental asociado a estos desechos, sino que también permite generar productos con valor agregado a partir de materiales considerados de bajo aprovechamiento.

De esta manera, el desarrollo de una bebida fermentada a partir de col y residuos de café representa una propuesta innovadora que integra beneficios probióticos y sostenibilidad, al combinar una fuente natural de microorganismos funcionales con el aprovechamiento de subproductos agroindustriales.

El objetivo de este estudio es analizar la valorización de la col y los residuos de café en el desarrollo de bebidas fermentadas probióticas de origen vegetal, a partir de la revisión y comparación de evidencia científica.  

\section{Materiales y métodos}

El presente estudio se desarrolló bajo un enfoque teórico-comparativo, con el objetivo de analizar la valorización de la col y los residuos de café en el desarrollo de bebidas fermentadas con potencial probiótico de origen vegetal. Para ello, se llevó a cabo una búsqueda sistemática de información científica utilizando palabras clave como “probiótico”, “vegetal”, “café” y “col”. 

A partir de esta búsqueda, se recopilaron y analizaron antecedentes relevantes, seleccionando aquellos estudios que abordaban el uso de matrices vegetales y residuos agroindustriales en procesos de fermentación con fines probióticos. Los artículos seleccionados fueron elegidos con base en su afinidad con el objetivo de la investigación y su pertinencia metodológica. 

Como parte del análisis, se tomó como referencia la metodología propuesta por \citet{guajardobarbosa2024}, la cual fue utilizada como modelo para comprender las etapas involucradas en la fermentación de matrices vegetales, incluyendo el acondicionamiento de la materia prima, el proceso de fermentación, el seguimiento de la cinética de crecimiento microbiano y la evaluación de metabolitos como el ácido láctico.

Posteriormente, se elaboró una matriz comparativa que permitió organizar la información recopilada y facilitar el análisis de los resultados reportados por diversos autores. Con base en esta matriz, se llevó a cabo un metaanálisis comparativo para evaluar tendencias en la viabilidad microbiana, producción de ácido láctico y aprovechamiento de residuos como los del café en el desarrollo de productos con potencial probiótico. 

Finalmente, los resultados obtenidos fueron sintetizados en un póster científico, en el que se integran los principales hallazgos relacionados con la valorización de la col y los residuos de café dentro del contexto de bebidas fermentadas de origen vegetal.

\section{Resultados}

La (...) presenta una síntesis de distintos artículos donde se realizaron fermentaciones con \textit{Lactobacillus plantarum} usando sustratos vegetales y diferentes condiciones experimentales, de manera general se demuestra como este microorganismos puede alcanzar concentraciones probióticas bajo condiciones experimentales controladas.

    %\input{tab:tabla_1}

La (...) presenta una comparación de diversos estudios que evalúan el uso de residuos de café como sustrato en procesos fermentativos con bacterias ácido-lácticas. En general, los resultados coinciden en que estos subproductos contienen compuestos fermentables que favorecen el crecimiento microbiano y la producción de metabolitos como el ácido láctico. Asimismo, se observa que la incorporación de residuos de café contribuye a mejorar la eficiencia del proceso fermentativo, evidenciada por una mayor actividad metabólica y un aprovechamiento más rápido de los sustratos. Estos hallazgos respaldan el potencial de los residuos de café como una alternativa viable para su valorización en sistemas fermentativos. 

    %\input{tab:tabla_2}
    
La Tabla 3 muestra que la media global de la concentración microbiana, expresada en logaritmo de UFC, es de 9.49 ± 0.63, lo que es equivalente  a 3 × 10⁹ UFC por mililitro, es decir, unos 3 mil millones de unidades formadoras de colonias por mL. El rango de valores reportados en los estudios va desde un mínimo de 7.02 log hasta un máximo de 13.01 log , superando en cada caso los umbrales mínimos propuestos por varios organismos como la FDA con valores mínimos de 10⁷ UFC , lo que demuestra  un nivel óptimo para productos comerciales.

    %\input{tab:tabla_3}
    
En la tabla 4 se muestra el error típico asociado a esta estimación es de 13.17, lo que refleja una precisión moderada. El intervalo de confianza del 95\% para el efecto combinado se extiende desde 4.12 hasta 68.57. Dado que este intervalo no incluye el cero (su límite inferior es 4.12, un valor positivo), se concluye que el efecto es estadísticamente significativo. Esta conclusión se confirma con el p-valor de 0.033, que al ser menor a 0.05 indica que la probabilidad de que este resultado se deba al azar es de solo un 3.3\%. Finalmente, la tabla reporta el rango de UFC reales obtenido en los estudios, que va desde 1.05 × 10⁷ hasta 1.02 × 10¹³ UFC por mL. Esto significa que en valores reales  existe una  amplia variabilidad en las magnitudes absolutas, aunque siempre con un impacto positivo y significativo de la fermentación. 

    %\input{tab:tabla_4}
    
La tabla 6 presenta los resultados de un análisis de meta-regresión realizado para evaluar si la unidad de medida (es decir, si la concentración de UFC se reportó en gramos  en mililitros) tiene un efecto significativo sobre el tamaño del efecto observado, aquí  el estadístico F es extremadamente bajo, con un valor de 8.76 × 10⁻⁵, lo que indica prácticamente nula variabilidad explicada por esta variable en los estudios incluidos, donde se puede concluir que no es una fuente significativa de heterogeneidad entre los estudios.
    
    %\input{tab:tabla_6}
    


\section{Discusión}

Los resultados del meta-análisis indican la presencia de un efecto combinado significativo  (p = 0.033), lo que sugiere que, en conjunto, las variables analizadas presentan una relación relevante dentro del proceso fermentativo estudiado. Sin embargo, el análisis también evidenció una heterogeneidad considerable entre los estudios (I$^{2}$ $\approx$ 79\%), lo que indica que existen diferencias importantes entre los datos incluidos. Esta variabilidad puede estar relacionada con factores como el tipo de sustrato utilizado, las condiciones de fermentación o las características de los microorganismos presentes.

Estos hallazgos son consistentes con lo reportado en la literatura sobre fermentación de vegetales. Diversos estudios han señalado que durante la fermentación de productos como la col se desarrolla una microbiota dominada por bacterias ácido-lácticas, principalmente de los géneros \textit{Lactobacillus} y \textit{Leuconostoc}, las cuales transforman los azúcares del vegetal en ácido láctico y generan un ambiente ácido que inhibe microorganismos indeseables \citep{dicagno2013, touret2018}. Este proceso no solo contribuye a la conservación del alimento, sino que también favorece la producción de compuestos bioactivos y el desarrollo de microorganismos con potencial probiótico.

La relevancia de estos microorganismos también ha sido respaldada por estudios experimentales sobre fermentaciones de vegetales. Por ejemplo, investigaciones sobre fermentaciones de chucrut han identificado diferentes cepas de bacterias ácido-lácticas con características asociadas al potencial probiótico, incluyendo tolerancia a condiciones ácidas y capacidad de sobrevivir al tránsito gastrointestinal \citep{semedolemsaddek2018}. Estas características coinciden con lo descrito en el marco teórico respecto a la capacidad de especies como \textit{Lactobacillus plantarum} para resistir condiciones adversas y contribuir al equilibrio de la microbiota intestinal.

Además, los resultados obtenidos pueden compararse con investigaciones de síntesis científica. Un meta-análisis sobre el efecto de alimentos fermentados y probióticos en trastornos gastrointestinales encontró que el consumo de productos fermentados que contienen bacterias ácido-lácticas puede generar mejoras significativas en la salud digestiva y en la función intestinal \citep{dimidi2014, eales2017}. Estos resultados sugieren que los microorganismos presentes en alimentos fermentados pueden tener efectos beneficiosos medibles en diferentes contextos clínicos. En este sentido, el efecto significativo observado en el presente meta-análisis podría interpretarse como evidencia adicional del papel funcional que desempeñan los procesos fermentativos y las bacterias ácido-lácticas en la generación de compuestos o condiciones favorables para la salud intestinal.

No obstante, el análisis de meta-regresión realizado mostró que la unidad variable (mL) no tuvo un efecto significativo (p = 0.993) sobre la magnitud del efecto. Esto sugiere que la variación en el volumen considerado en los estudios no explica las diferencias observadas entre los resultados. Este comportamiento puede deberse a que otros factores, como la composición microbiana inicial, la concentración de sal, la disponibilidad de azúcares o las condiciones ambientales de fermentación, tienen una influencia más importante sobre la dinámica del proceso fermentativo.

Desde el punto de vista aplicado, estos resultados refuerzan el interés en el desarrollo de alimentos y bebidas fermentadas con potencial funcional, especialmente aquellos que contienen bacterias ácido-lácticas. Diversos autores han señalado que los alimentos fermentados pueden contribuir a mejorar la diversidad de la microbiota intestinal, favorecer la digestión y promover la producción de metabolitos beneficiosos para el organismo \citep{marco2017}. El uso de sustratos alternativos como residuos agroindustriales representa una estrategia prometedora para el desarrollo de nuevos productos biotecnológicos con valor agregado.

En este sentido, la integración de ambos componentes como la col como fuente de microorganismos probióticos y los residuos de café como sustrato prebiótico representan una estrategia innovadora que podría dar lugar a sistemas simbióticos, este tipo de sistemas tiene el potencial de mejorar la supervivencia y actividad de los microorganismos durante la fermentación y posteriormente en el tracto gastrointestinal \citep{suez2019}.

A pesar de estos hallazgos, es importante considerar algunas limitaciones del presente análisis. En primer lugar, el número de estudios incluidos fue relativamente reducido, lo cual puede influir en la precisión de las estimaciones obtenidas. Asimismo, la elevada heterogeneidad observada sugiere que los estudios analizados presentan diferencias metodológicas importantes, lo que dificulta la comparación directa de los resultados. Estas variaciones pueden estar relacionadas con factores como el tipo de sustrato vegetal utilizado, las condiciones de fermentación, las cepas microbianas presentes o las metodologías empleadas para evaluar los efectos.

Por lo tanto, futuras investigaciones deberían centrarse en estandarizar las condiciones experimentales en estudios de fermentación y en ampliar el número de investigaciones disponibles para meta-análisis. También sería relevante explorar con mayor profundidad la interacción entre diferentes sustratos vegetales, microorganismos fermentativos y compuestos bioactivos generados durante el proceso. De igual manera, se requieren más estudios que evalúen los efectos de alimentos fermentados en la salud humana mediante diseños experimentales controlados, lo que permitiría fortalecer la evidencia científica sobre sus beneficios potenciales. 

\section{Conclusión}

A partir de la revisión y comparación de la evidencia científica, se puede concluir que la col y los residuos de café presentan un potencial importante para el desarrollo de bebidas fermentadas probióticas de origen vegetal. La col destaca como una matriz adecuada para la fermentación láctica debido a su contenido de azúcares fermentables y a la presencia natural de bacterias ácido-lácticas como \textit{Lactobacillus plantarum}, las cuales han demostrado alcanzar concentraciones compatibles con efectos probióticos en distintos estudios.

Por otro lado, los residuos de café representan una alternativa interesante para su aprovechamiento biotecnológico, ya que contienen compuestos bioactivos y oligosacáridos que pueden actuar como sustratos o compuestos con posible efecto prebiótico, favoreciendo la actividad metabólica de bacterias benéficas durante los procesos fermentativos. La evidencia analizada sugiere que su incorporación puede contribuir a mejorar la dinámica de fermentación y al mismo tiempo promover la valorización de subproductos agroindustriales.

Combinar estos dos componentes podría favorecer el desarrollo de sistemas simbióticos. Abriendo la posibilidad de desarrollar nuevas bebidas funcionales de orígen vegetal, así como también brinda una alternativa más sostenible dentro del aprovechamiento de recursos alimenticios.

Sin embargo, cabe mencionar que aún es necesario realizar más investigaciones experimentales para así evaluar ciertos aspectos, como la estandarización del proceso, estabilidad microbiológica, efectos a la salud humana, etcétera. Para así determinar con mayor precisión la viabilidad de este tipo de bebidas en el sector alimentario. 


% ------------------------------------------------------ %

    \bibliographystyle{apalike-ejor}
    \bibliography{art_02/bib}

% ------------------------------------------------------ %


\end{multicols}